\slajdid{09_1-s005}
\begin{frame}[label=09_1-s005]
\gusttekst
Ovi izrazi, odnosno model je pokazivao dobre rezultate kada su dimenzije tranzistora bile „velike“ reda $\mu m$. Međutim prilikom eksperimenata sa tranzistorima čije su dimenzije manje i reda $nm$ pojavio se efekat, gledajući izlazne karakteristike, da tranzistor ulazi „ranije“ u oblast zasićenja odnosno pri manjim naponima $V_{DS}$.\par\medskip
Uočeno je da pri „velikim“ poljima u kanalu dolazi do zasićenja brzine nosilaca, što efektivno ima uticaj da tranzistor „ranije“ uđe u režim zasićenja.\par\medskip
\centering Uvedeni pojmovi dugačak i kratak kanal\par
\begin{columns}[c,onlytextwidth]\column{.40\textwidth}\centering
\tikz{\draw[->,DEplava](1,0)--(0,-.6);\draw[->,DEplava](2.8,0)--(3.3,-.6);}\par
\begin{tikzpicture}[x=.65cm,y=.65cm,font=\sitneoznake,>=Latex]\draw[->](-.15,0)--(2.5,0)node[below]{$E$};\draw[->](0,-.15)--(0,2.3)node[left]{$v$};\draw[<->](.85,0)arc[start angle=0,end angle=50,radius=.85];\node at(.35,.17){$\mu$};\draw(0,0)--(1.45,1.8);\node[align=center]at(1.35,2.4){dugačak\\kanal};\end{tikzpicture}
\hfill\begin{tikzpicture}[x=.65cm,y=.65cm,font=\sitneoznake,>=Latex]\draw[->](-.15,0)--(2.5,0)node[below]{$E$};\draw[->](0,-.15)--(0,2.3)node[left]{$v$};\draw[<->](.85,0)arc[start angle=0,end angle=50,radius=.85];\node at(.35,.17){$\mu$};\draw(0,0)..controls(.7,.9)and(.9,1.05)..(1.1,1.1)..controls(1.5,1.2)and(1.9,1.2)..(2.3,1.2);\draw[dashed](0,0)--(1.2,1.6);\draw[dashed](0,1.2)node[left]{$v_{sat}$}--(2.3,1.2);\draw[dashed](1.1,0)node[below]{$E_C$}--(1.1,1.6);\node[align=center]at(1.5,2.4){kratak\\kanal};\end{tikzpicture}

\column{.58\textwidth}
\[v_n=\mu_n\frac E{1+\frac E{E_{Cn}}}\quad za\ E\leq E_{Cn}\]
\[v_n=v_{nsat}\quad za\ E\geq E_{Cn}\]
\raggedright gde je E polje koje deluje na nosioce, $E_{Cn}$ kritično polje pri kojem nastaje efekat zasićenja brzine nosilaca a $v_{nsat}=\frac{\mu_nE_{Cn}}2$.
\end{columns}
\end{frame}
